% £robl€m 002: A bushel of apples
% Body only -- no \documentclass, no preamble packages. Title, dates, status live in the
% front matter of _posts/2026-11-04-problem-002-a-buchel-of-apples.html.
% Macros defined here are expanded by pandoc:
\newcommand{\vect}[1]{\boldsymbol{#1}}

Turning a classroom into a lab.

\section*{Setting}
Your university/school has equipped rooms with multiple (apple) desktop computers which are unused when there are
no classes. During this idle time, these clusters of c/gpus shall be used for scientific tasks. Given the
limitations placed by the start of the next class, the relatively slow connection between the computing nodes,
and all issues related to the closedness of iOS, realize a framework which would allow me to distribute $\mathcal{O}(10^4)$ runs of
a single sequence of codes (FORTRAN, python, &c.) with typical runtime $\tau=\mathcal{O}(10m)$ onto the classroom.
The framework should be controllable from my Linux terminal, and the results of each client run should be stored there temporarily
before sending them to me.

If you plan to engage in this problem, contact me for sample code.

\section*{Deliverable}
\begin{itemize}
  \item Flow diagram of the computation starting from its trigger on the Linux terminal to the final result collection
  \item practical implementation respecting the boundary conditions, i.e., runtime estimate and possible job rejection
\end{itemize}
